%========================= METODOLOGIA =========================
\section{Metodologia}

% RASCUNHO gerado a partir do código e da documentação do protótipo.
% Revisar, ajustar os números que ainda dependem da coleta e inserir
% as citações dos artigos selecionados na revisão.

Este trabalho caracteriza-se como uma pesquisa aplicada, de natureza
exploratória e abordagem predominantemente quantitativa, organizada em quatro
etapas: revisão da literatura, desenvolvimento do artefato, coleta de dados com
participantes e análise dos resultados.

\subsection{Revisão da literatura}

A fundamentação teórica foi construída a partir de uma busca sistematizada nas
bases Scopus e SpringerLink, restrita a publicações dos últimos cinco anos. A
expressão de busca combinou quatro blocos de termos --- tecnologia, condição,
estímulo e procedimento --- conforme apresentado a seguir:

\begin{quote}\small\singlespacing
\texttt{( "virtual reality" OR "immersive environment" OR WebXR )}\\
\texttt{AND ( phobia OR "specific phobia" OR fear OR anxiety )}\\
\texttt{AND ( spider OR snake OR cockroach OR insect OR animal}\\
\texttt{\phantom{AND (} OR arachnophobia OR ophidiophobia OR entomophobia )}\\
\texttt{AND ( exposure OR "exposure therapy" OR desensitization OR habituation )}
\end{quote}

% TODO: preencher depois de rodar a busca
Os resultados foram filtrados por leitura de título e resumo, chegando-se a um
conjunto final de \textbf{[N]} artigos, que fundamentam as decisões de projeto
descritas na Seção~\ref{sec:desenvolvimento} e a interpretação dos resultados.

\subsection{Participantes}

% TODO: confirmar números após a coleta
A validação foi conduzida com estudantes dos cursos de Ciência da Computação e
de Psicologia da Faculdade Nova Roma. A participação foi voluntária e mediante
assinatura do Termo de Consentimento Livre e Esclarecido (Apêndice~A).

Cabe registrar que \textbf{o sistema não constitui tratamento psicológico ou
médico e não realiza diagnóstico}. A avaliação conduzida restringe-se a
usabilidade, imersão, interação e percepção do estímulo, não se propondo a medir
desfecho clínico de qualquer natureza.

\subsection{Procedimento}

Cada sessão seguiu o mesmo roteiro:

\begin{enumerate}
  \item apresentação do estudo e leitura do termo de consentimento;
  \item cadastro na plataforma, com registro de e-mail, sexo e data de
  nascimento;
  \item questionário inicial de perfil, no qual o participante declara, em
  escala de 0 a 10, o grau de medo que atribui a cada um dos cinco animais
  contemplados;
  \item escolha do estímulo e execução da experiência, percorrendo os níveis de
  exposição na ordem crescente, com liberdade para interromper a qualquer
  momento;
  \item questionário final de percepção.
\end{enumerate}

\subsection{Variáveis coletadas}

Os dados foram registrados automaticamente pela aplicação durante a experiência
e associados a um código de participante (P01, P02, \ldots), de modo que a
análise é conduzida sobre dados anonimizados. O Quadro a seguir relaciona cada
variável comportamental à sua forma de medida.

\begin{table}[htbp]
\centering
\caption{Variáveis comportamentais registradas durante a experiência.}
\label{tab:variaveis}
\begin{tabular}{ll}
\toprule
\textbf{Variável} & \textbf{Como é medida} \\
\midrule
Aproximação          & menor distância alcançada até o estímulo, em metros \\
Afastamento          & número de recuos e distância média mantida \\
Agitação             & variação angular da câmera por segundo \\
Atenção ao estímulo  & proporção do tempo com o estímulo no campo de visão \\
Exploração           & deslocamento acumulado e animais encontrados \\
Interação            & ocorrência de toque no estímulo (Nível 5) \\
Progressão           & nível máximo alcançado e tempo por nível \\
Desconforto          & autorrelato em escala de 0 a 10, por nível \\
\bottomrule
\end{tabular}
\end{table}

\subsection{Tratamento dos dados}

Os registros são gravados em banco de dados relacional e exportados para
análise descritiva. As figuras e tabelas apresentadas na Seção de resultados são
geradas por scripts versionados em \texttt{data/analytics}, a partir dos dados
já anonimizados, o que torna a análise reprodutível sem expor informação
identificável dos participantes.
